% from ega1-4-fr.tex
\oldpage[I]{122}
de préschémas (\hyperref[I.2.2.1-fr]{2.2.1}), qui, en vertu de ce qui précède,
est tel que $f=j\circ g$, d'où la proposition.

\begin{corollary}[4.1.10]
\label{I.4.1.10-fr}
Pour qu'un morphisme d'injection $Z\to X$ soit \emph{majoré par le morphisme
d'injection} $Y\to X$, il faut et il suffit que $Z$ soit un sous-préschéma de
$Y$.
\end{corollary}

On écrit alors $Z\leq Y$, et cette relation est évidemment une
\emph{relation d'ordre} dans l'ensemble des sous-préschémas de $X$.

\subsection{Morphismes d'immersion}
\label{subsection:I.4.2-fr}

\begin{definition}[4.2.1]
\label{I.4.2.1-fr}
On dit qu'un morphisme $f:Y\to X$ est une \emph{immersion} (resp. une
\emph{immersion fermée}, une \emph{immersion ouverte}) s'il se factorise en
$Y\xrightarrow{g}Z\xrightarrow{j}X$, où $g$ est un isomorphisme, $Z$ un
sous-préschéma de $X$ (resp. un sous-préschéma fermé, un sous-préschéma induit
sur un ouvert) et $j$ le morphisme d'injection.
\end{definition}

Le sous-préschéma $Z$ et l'isomorphisme $g$ sont alors déterminés de façon
\emph{unique}, car si $Z'$ est un second sous-préschéma de $X$, $j'$ l'injection
$Z'\to X$ et $g'$ un isomorphisme $Y\to Z'$ tel que
$j\circ g=j'\circ g'$, on en déduit $j'=j\circ g\circ {g'}^{-1}$, d'où
$Z'\leq Z$ (\hyperref[I.4.1.10-fr]{4.1.10}), et on montre de même que
$Z\leq Z'$, donc $Z'=Z$, et comme $j$ est un monomorphisme de préschémas,
$g'=g$.

On dit que $f=j\circ g$ est la \emph{factorisation canonique} de l'immersion
$f$, et le sous-préschéma $Z$ et l'isomorphisme $g$ sont dits
\emph{associés à $f$}.

Il est clair qu'une immersion est un \emph{monomorphisme} de préschémas
(\hyperref[I.4.1.7-fr]{4.1.7}) et \emph{a fortiori} un morphisme
\emph{radiciel} (\hyperref[I.3.5.4-fr]{3.5.4}).

\begin{proposition}[4.2.2]
\label{I.4.2.2-fr}
\begin{enumerate}
  \item[a)] Pour qu'un morphisme $f=(\psi,\theta):Y\to X$ soit une
  \emph{immersion ouverte}, il faut et il suffit que $\psi$ soit un
  homéomorphisme de $Y$ sur une partie ouverte de $X$, et que pour tout
  $y\in Y$, l'homomorphisme
  $\theta_y^\sharp:\mathscr{O}_{\psi(y)}\to\mathscr{O}_y$ soit bijectif.
  \item[b)] Pour qu'un morphisme $f=(\psi,\theta):Y\to X$ soit une
  \emph{immersion} (resp. une \emph{immersion fermée}), il faut et il suffit
  que $\psi$ soit un homéomorphisme de $Y$ sur une partie localement fermée
  (resp. fermée) de $X$, et que pour tout $y\in Y$, l'homomorphisme
  $\theta_y^\sharp:\mathscr{O}_{\psi(y)}\to\mathscr{O}_y$ soit surjectif.
\end{enumerate}
\end{proposition}

\begin{enumerate}
  \item[a)] Les conditions sont évidemment nécessaires. Inversement, si elles
  sont remplies, il est clair que $\theta^\sharp$ est un isomorphisme de
  $\mathscr{O}_Y$ sur $\psi^*(\mathscr{O}_X)$, et
  $\psi^*(\mathscr{O}_X)$ est le faisceau déduit par transport de structure au
  moyen de $\psi^{-1}$ à partir de $\mathscr{O}_X|\psi(Y)$~; d'où la
  conclusion.
  \item[b)] Les conditions étant évidemment nécessaires, prouvons qu'elles
  sont suffisantes. Considérons d'abord le cas particulier où l'on suppose que
  $X$ est un schéma affine et que $Z=\psi(Y)$ est \emph{fermé} dans $X$. On
  sait (\hyperref[0.3.4.6-fr]{0, 3.4.6}) que $\psi_*(\mathscr{O}_Y)$ a alors
  pour support $Z$ et que si l'on désigne par $\mathscr{O}'_Z$ sa restriction
  à $Z$, l'espace annelé $(Z,\mathscr{O}'_Z)$ se déduit de
  $(Y,\mathscr{O}_Y)$ par transport de structure au moyen de
  l'homéomorphisme $\psi$, considéré comme application de $Y$ sur $Z$.
  Montrons que $f_*(\mathscr{O}_Y)=\psi_*(\mathscr{O}_Y)$ est un
  $\mathscr{O}_X$-Module \emph{quasi-cohérent}. En effet, pour tout
  $x\notin Z$, $\psi_*(\mathscr{O}_Y)$, restreint à un voisinage convenable de
  $x$, est nul. Si au contraire $x\in Z$, on a $x=\psi(y)$ pour un $y\in Y$
  bien déterminé~; soit $V$ un voisinage ouvert affine de $y$ dans $Y$~;
  $\psi(V)$ est alors ouvert dans $Z$, donc trace sur $Z$ d'un ouvert $U$ de
  $X$, et la restriction à $U$ de $\psi_*(\mathscr{O}_Y)$ est identique à la
  restriction à $U$ de l'image
